\section{\texorpdfstring{第~\ref{chap:geometric-mechanics}~章}{几何力学}}\label{sec:solutions-geometric-mechanics}
\subsection{\texorpdfstring{第~\ref{sec:mechanical-systems}~节}{力学系统}}\label{subsec:solutions-mechanical-systems}

\begin{description}[font=\normalfont,leftmargin=2em,labelwidth=1.5em]
\RGsolution{ex:prove-levi-civita-under-conformal-change}{levi-civita-under-conformal-change}
由 Koszul 公式,
\begin{align*}
  2\langle\langle\widetilde\nabla_XY,Z\rangle\rangle
  &=X\cdot\langle\langle Y,Z\rangle\rangle+Y\cdot\langle\langle X,Z\rangle\rangle
     -Z\cdot\langle\langle X,Y\rangle\rangle\\
  &\quad-\langle\langle[X,Z],Y\rangle\rangle-\langle\langle[Y,Z],X\rangle\rangle
     +\langle\langle[X,Y],Z\rangle\rangle.
\end{align*}
例如, 注意
\begin{align*}
  X\cdot\langle\langle Y,Z\rangle\rangle
  &=X\cdot\left(e^{2\rho}\langle Y,Z\rangle\right)\\
  &=e^{2\rho}X\cdot\langle Y,Z\rangle+2\,\mathrm{d}\rho(X)e^{2\rho}\langle Y,Z\rangle,
\end{align*}
则显然有
\begin{align*}
  2\langle\langle\widetilde\nabla_XY,Z\rangle\rangle
  &=2e^{2\rho}\langle\nabla_XY,Z\rangle+2\,\mathrm{d}\rho(X)e^{2\rho}\langle Y,Z\rangle\\
  &\quad+2\,\mathrm{d}\rho(Y)e^{2\rho}\langle X,Z\rangle-2\,\mathrm{d}\rho(Z)e^{2\rho}\langle X,Y\rangle\\
  &=2\langle\langle\nabla_XY,Z\rangle\rangle+2\langle\langle\mathrm{d}\rho(X)Y,Z\rangle\rangle\\
  &\quad+2\langle\langle\mathrm{d}\rho(Y)X,Z\rangle\rangle
       -2e^{2\rho}\langle\operatorname{grad}\rho,Z\rangle\langle X,Y\rangle.
\end{align*}
由于
\[
  e^{2\rho}\langle\operatorname{grad}\rho,Z\rangle
    =\langle\langle\operatorname{grad}\rho,Z\rangle\rangle
\]
且 $Z$ 是任取的, 得
\[
  \widetilde\nabla_XY=\nabla_XY+\mathrm{d}\rho(X)Y+\mathrm{d}\rho(Y)X
     -\langle X,Y\rangle\operatorname{grad}\rho.
\]

\RGsolution{ex:prove-criterion-for-reparameterized-geodesic}{criterion-for-reparameterized-geodesic}
设 $s:I\to J$ 是微分同胚, $\gamma:J\to M$ 是 $c:I\to M$ 的重新参数化, 由
\[
  c(t):=\gamma(s(t))
\]
定义. 有
\setcounter{equation}{28}
\begin{equation}\label{eq:reparameterized-curve-chain-rule}
  \dot c(t)=\dot\gamma(s(t))\frac{\mathrm{d}s}{\mathrm{d}t}(t),
\end{equation}
从而
\begin{align*}
  \nabla_{\dot c}\dot c
  &=\nabla_{\dot c}\left(\frac{\mathrm{d}s}{\mathrm{d}t}\dot\gamma\right)
    =\left(\frac{\mathrm{d}^2s}{\mathrm{d}t^2}\right)\dot\gamma
      +\left(\frac{\mathrm{d}s}{\mathrm{d}t}\right)
          \nabla_{\frac{\mathrm{d}s}{\mathrm{d}t}\dot\gamma}\dot\gamma\\
  &=\frac{\mathrm{d}}{\mathrm{d}t}\left(\log\left|\frac{\mathrm{d}s}{\mathrm{d}t}\right|\right)\dot c
      +\left(\frac{\mathrm{d}s}{\mathrm{d}t}\right)^2\nabla_{\dot\gamma}\dot\gamma.
\end{align*}
因此, 若 $c$ 是测地线 $\gamma$ 的重新参数化, 则 $c$ 满足
\[
  \nabla_{\dot c}\dot c=f(t)\dot c,
\]
其中
\[
  f(t)=\frac{\mathrm{d}}{\mathrm{d}t}\left(\log\left|\frac{\mathrm{d}s}{\mathrm{d}t}\right|\right).
\]
另一方面, 若 $c$ 满足 \eqref{eq:reparameterized-curve-chain-rule}, 则由
\[
  s(t)=\int\exp\left(\int f(t)\,\mathrm{d}t\right)\mathrm{d}t
\]
确定的 $c$ 的重新参数化 $\gamma$ 是一条测地线.

\RGsolution{ex:derive-newton-equation-in-coordinates}{derive-newton-equation-in-coordinates}
回顾, $TM$ 上的局部坐标 $(x^1,\ldots,x^n,v^1,\ldots,v^n)$ 参数化的向量
\[
  \sum_{i=1}^nv^i\frac{\partial}{\partial x^i}
\]
在坐标为 $(x^1,\ldots,x^n)$ 的点处与 $M$ 相切. 因此, 有
\[
  K(x^1,\ldots,x^n,v^1,\ldots,v^n)=\frac12\sum_{i,j=1}^n g_{ij}(x^1,\ldots,x^n)v^iv^j,
\]
其中
\[
  g_{ij}=\left\langle\frac{\partial}{\partial x^i},\frac{\partial}{\partial x^j}\right\rangle
\]
是此坐标系下的度量分量. 从而
\[
  \frac{\partial K}{\partial v^i}=\sum_{j=1}^n g_{ij}v^j,
\]
所以
\[
  \frac{\partial K}{\partial v^i}(x(t),\dot{x}(t))=\sum_{j=1}^n g_{ij}(x(t))\dot{x}^j(t),
\]
进而
\begin{align*}
  \frac{\mathrm{d}}{\mathrm{d}t}
   \left(\frac{\partial K}{\partial v^i}(x(t),\dot{x}(t))\right)
  &=\sum_{j=1}^n g_{ij}(x(t))\ddot{x}^j(t)\\
  &\quad+\sum_{j,k=1}^n\frac{\partial g_{ij}}{\partial x^k}(x(t))\dot{x}^k(t)\dot{x}^j(t).
\end{align*}
此外,
\[
  \frac{\partial K}{\partial x^i}=\frac12\sum_{j,k=1}^n
   \frac{\partial g_{jk}}{\partial x^i}v^jv^k,
\]
所以
\[
  \frac{\partial K}{\partial x^i}(x(t),\dot{x}(t))
    =\frac12\sum_{j,k=1}^n\frac{\partial g_{jk}}{\partial x^i}(x(t))\dot{x}^j(t)\dot{x}^k(t).
\]
因此
\begin{align*}
  &\frac{\mathrm{d}}{\mathrm{d}t}
   \left(\frac{\partial K}{\partial v^i}(x(t),\dot{x}(t))\right)
   -\frac{\partial K}{\partial x^i}(x(t),\dot{x}(t))\\
  &\qquad=\sum_{j=1}^n g_{ij}(x(t))\ddot{x}^j(t)
    +\sum_{j,k=1}^n\left(\frac{\partial g_{ij}}{\partial x^k}(x(t))
       -\frac12\frac{\partial g_{jk}}{\partial x^i}(x(t))\right)\dot{x}^j(t)\dot{x}^k(t).
\end{align*}
另一方面, 若 $v,w\in T_pM$ 写为
\[
  v=\sum_{i=1}^nv^i\frac{\partial}{\partial x^i},\quad
  w=\sum_{i=1}^nw^i\frac{\partial}{\partial x^i},
\]
则
\[
  \mu(v)(w)=\sum_{i,j=1}^n g_{ij}v^i w^j
      =\sum_{i,j=1}^n g_{ij}v^i\,\mathrm{d}x^j(w),
\]
因而
\[
  \mu(v)=\sum_{i,j=1}^n g_{ij}v^i\,\mathrm{d}x^j
     =\sum_{i,j=1}^n g_{ij}v^j\,\mathrm{d}x^i.
\]
因此
\[
  \mu\left(\frac{D\dot c}{\mathrm{d}t}\right)(t)
  =\sum_{i,j=1}^n g_{ij}(x(t))
   \left(\ddot{x}^j(t)+\sum_{k,l=1}^n\Gamma^j_{kl}(x(t))\dot{x}^k(t)\dot{x}^l(t)\right)\mathrm{d}x^i.
\]
由于
\begin{align*}
  \sum_{j=1}^n g_{ij}\Gamma^j_{kl}
  &=\frac12\sum_{j,m=1}^n g_{ij}g^{jm}
     \left(\frac{\partial g_{ml}}{\partial x^k}+\frac{\partial g_{mk}}{\partial x^l}
         -\frac{\partial g_{kl}}{\partial x^m}\right)\\
  &=\frac12\left(\frac{\partial g_{il}}{\partial x^k}+\frac{\partial g_{ik}}{\partial x^l}
         -\frac{\partial g_{kl}}{\partial x^i}\right),
\end{align*}
所以
\begin{align*}
  &\sum_{j,k,l=1}^n g_{ij}(x(t))\Gamma^j_{kl}(x(t))\dot{x}^k(t)\dot{x}^l(t)\\
  &\quad=\frac12\sum_{k,l=1}^n
    \left(\frac{\partial g_{il}}{\partial x^k}(x(t))
              +\frac{\partial g_{ik}}{\partial x^l}(x(t))
              -\frac{\partial g_{kl}}{\partial x^i}(x(t))\right)\dot{x}^k(t)\dot{x}^l(t)\\
  &\quad=\sum_{j,k=1}^n\left(\frac{\partial g_{ij}}{\partial x^k}(x(t))
            -\frac12\frac{\partial g_{jk}}{\partial x^i}(x(t))\right)\dot{x}^j(t)\dot{x}^k(t),
\end{align*}
证明完毕.
\end{description}
